% =============================================================================
%  Breadcrumbs — tables to accompany Figure 4.
%
%  Paste into main.tex near Section 7.3. Needs booktabs.
%  Regenerate the numbers with:
%      cd "Breadcrumbs Project" && python -m model.run train --rounds 15 --seeds 5
%  Source of truth: Breadcrumbs Project/model/artefacts/training.json
% =============================================================================
%  NOTE ON DUPLICATION
%    main.tex carries an inlined copy of this content, because the paper draws
%    every figure in TikZ so that nothing has to be uploaded alongside it. This
%    standalone file exists for previewing one figure without compiling the whole
%    report. If you change the numbers, change both, or the preview will disagree
%    with the paper. Regenerate from model/artefacts/training.json.

\begin{table}[htbp]
\centering
\caption{Federated continual learning on the document corpus, mean and standard
deviation over five seeds. Six sites act as federated clients with a
Dirichlet($\alpha{=}0.6$) partition; the three violation waves are the continual
stages. Balanced accuracy is reported because 4.11\% of the corpus is anomalous
and plain accuracy would reward a detector that finds nothing; chance is 50.0.
Recall is measured as filing a document under the correct wave, not merely
flagging it. Every figure is measured at the 10\% false-positive budget of
Table~\ref{tab:corpus_operating}. Backward transfer is $-25.2\,\pm\,3.5$ for
sequential averaging and $-12.7\,\pm\,2.0$ for prototype replay. Unlike
Table~\ref{tab:results}, these are row-level documents with nine anomaly kinds,
not Gaussian clusters.}
\label{tab:corpus}
\small
\setlength{\tabcolsep}{4.5pt}
\begin{tabular}{@{}l c c c c c@{}}
\toprule
\textbf{Method} & \textbf{Wave 1} & \textbf{Wave 2} & \textbf{Wave 3} &
\textbf{Average} & \textbf{Forgetting} \\
\midrule
Chance & 50.0 & 50.0 & 50.0 & 50.0 & n/a \\
Sequential federated averaging & 50.0\,$\pm$\,0.1 & 92.4\,$\pm$\,6.7 &
  \textbf{68.8\,$\pm$\,0.9} & 70.4\,$\pm$\,2.1 & 16.8\,$\pm$\,2.3 \\
\textbf{Prototype replay (ours)} & \textbf{67.8\,$\pm$\,3.4} &
  \textbf{99.7\,$\pm$\,0.5} & 63.4\,$\pm$\,2.4 & \textbf{77.0\,$\pm$\,1.5} &
  \textbf{8.5\,$\pm$\,1.3} \\
\midrule
\textit{Centralised pooling (not permitted)} & \textit{85.7\,$\pm$\,2.7} &
  \textit{100.0\,$\pm$\,0.0} & \textit{63.9\,$\pm$\,2.0} & \textit{83.2\,$\pm$\,0.3} &
  \textit{n/a} \\
\bottomrule
\end{tabular}
\end{table}

\begin{table}[htbp]
\centering
\caption{Choosing the operating point. A four-way $\arg\max$ has no adjustable
threshold, and the point it happens to land on gave a false-positive rate of
39.1\,$\pm$\,8.0, which is unusable for an auditor and unstable across seeds. Scoring
instead on $P(\text{anomalous})$ against a threshold chosen on a held-out 15\%
split of the \emph{training} shards makes the false-positive rate a quantity the
consortium sets. The benchmarks are never used to choose it. ROC-AUC is
$0.894\,\pm\,0.004$ and does not depend on the threshold at all. Note that the
5\% budget dominates the 10\% one: it scores higher balanced accuracy while
flagging half as many clean documents.}
\label{tab:corpus_operating}
\small
\begin{tabular}{@{}l r r r@{}}
\toprule
\textbf{False-positive budget} & \textbf{Detection} & \textbf{Realised FP} &
\textbf{Balanced} \\
\midrule
1\%  & 67.3\,$\pm$\,1.9 &  1.0\,$\pm$\,0.2 & 83.1\,$\pm$\,1.0 \\
5\%  & 74.4\,$\pm$\,1.8 &  5.1\,$\pm$\,0.5 & \textbf{84.7\,$\pm$\,0.9} \\
10\% & 77.3\,$\pm$\,2.3 & 10.4\,$\pm$\,0.8 & 83.5\,$\pm$\,1.2 \\
\midrule
\textit{$\arg\max$, no threshold (previous)} & \textit{88.7} & \textit{43.2} &
  \textit{72.8} \\
\bottomrule
\end{tabular}
\end{table}

\begin{table}[htbp]
\centering
\caption{Detection by anomaly kind on the held-out cross-wave benchmark, for the
prototype-replay detector after all three waves, at the 10\% budget. Mean and
standard deviation over the same five seeds. The false-positive rate is given in
the same table rather than a footnote, because a per-kind detection rate without
it is not a result: a model that answers ``anomalous'' more often scores better
on every row and is worse at the job. \texttt{cross\_inconsistency} is the
control. Those documents are internally consistent and contradict a
\emph{different} document of the same period, so no single-document feature can
see them, and at 8.3\% they are flagged \emph{less} often than clean documents
are, which is the correct result and the one we would have had to explain away
had it come out otherwise.}
\label{tab:corpus_kinds}
\small
\begin{tabular}{@{}l r r@{}}
\toprule
\textbf{Anomaly kind} & \textbf{Detected} & \textbf{$n$} \\
\midrule
Back-dated signature          & 100.0\,$\pm$\,0.0 &  98 \\
Benford first-digit violation & 100.0\,$\pm$\,0.0 & 145 \\
Certificate / CAS check digit & 100.0\,$\pm$\,0.0 & 125 \\
Arithmetic inconsistency      &  99.2\,$\pm$\,0.4 &  95 \\
Duplicated reference          &  90.4\,$\pm$\,13.5 & 160 \\
Implausible roundness         &  80.0\,$\pm$\,15.5 &  39 \\
Overtime above the ceiling    &  70.9\,$\pm$\,17.6 &  11 \\
Site-relative outlier         &  50.1\,$\pm$\,1.6 & 284 \\
\midrule
\textit{Clean documents (false positives)} & \textit{10.4\,$\pm$\,0.8} & \textit{971} \\
\textit{Cross-document inconsistency (control)} & \textit{8.3\,$\pm$\,5.1} & \textit{70} \\
\bottomrule
\end{tabular}
\end{table}
